% page 139 du fac-simile (image p-138 du scan)
% bloc 165.1 x 246.3 mm ; marge gauche 36.9 mm
\pgc{15.240mm}{0.000mm}{
\l{}
\l{}
\l{\cel{55}129}
\l{}
\l{}
\l{\sou{elefanto}.-(\sou{zool.}) Mamifero de la klaso \textquotedbl{}pakidermi\textquotedbl{}, remarkinda}
\l{\cel{3}pro la maso enorma di lua korpo, lua pelo sen-pila e rugoza,}
\l{\cel{3}la dentegi qui garnisas lua maxilo, e lua nazo longigita a}
\l{\cel{3}rostro. - L.\cel{1}\sou{elephas}. - DEFIRS.}
\l{}
\l{\sou{elefantiazo}.- (\sou{patol.}) Quaza letro, karakterizata da tuberkuli}
\l{\cel{3}di la pelo qui igas olca rugoza quale ta di la elefanto. -}
\l{\cel{3}DEFIRS.}
\l{}
\l{\sou{eleganta}.- I. Karakterizata da distingateso plena ye gracio e}
\l{\cel{3}ye sen-jeneso. - II. Di qua la manieri esas gracioza e sen-}
\l{\cel{3}pene tala. - DEFIRS.}
\l{}
\l{\sou{elegio}.- (en\cel{1}\sou{la epoki antiqua}).- Mikra poemo di kataktero}
\l{\cel{3}melankolioza e tenera. - DEFIRS.}
\l{}
\l{\sou{elektar}. (\sou{trans., por...-eso}) (\sou{aludante plura personi}).}
\l{\cel{3}Selektar per voto (por ofico, plaso, posteno). - EFIS.}
\l{}
\l{\sou{elektro}.- (\sou{Fiziko}). Proprajo quan ula korpi posedas od aquiras,}
\l{\cel{3}kande li frotesas, kompresesas, e c., atraktar o retropulsar}
\l{\cel{3}altra korpi, produktar cintili, deskompozar kemiale, sukusar}
\l{\cel{3}la nervi, e c.-\cel{2}DEFIRS.}
\l{}
\l{\sou{elektrodo}. (\sou{Fiziko)}.\cel{2}Konduktivo quan onu igas komunikar kun}
\l{\cel{3}calatere la baterio, talatere medio a qua la fluo elektra}
\l{\cel{3}efikas kemiale. - DEFIRS.}
\l{}
\l{\sou{elektroforo}. - Aparato qua uzesas por akumular elektro. - DEFIRS.}
\l{}
\l{\sou{elektrolito}. Korpo qua, liquida, esas deskompozebla da la}
\l{\cel{3}elektro-fluo. - DEFIRS.}
\l{}
\l{\sou{elektrolizar}. (\sou{trans.})\cel{2}Deskompozar kemiale per la fluo elektra.}
\l{\cel{3}- DEFIRS.}
\l{}
\l{\sou{elektrologio}.\cel{2}Parto di fiziko qua traktas la fenomeni elektrala}
\l{\cel{3}e magnetala (elektro-magnetala), e lia relati kun altra feno-}
\l{\cel{3}meni. - DEFIRS.}
\l{}
\l{\sou{elektrometro}. Instrumento uzata por mezurar la intenseso o la}
\l{\cel{3}naturo (pozitiva o negativa) di la fluo elektra ye qua karg-}
\l{\cel{3}esas korpo. - DEFIRS.}
\l{}
\l{\sou{elektrometrio}. Ensemblo de la metodi pri la mezuro di la elektro-}
\l{\cel{3}-grandori (intenseso, rezistiveso, e c.) - DEFIRS.}
\l{}
\l{\sou{elektromobilo}. Veturo movata per elektro-motoro (vice per kavali}
\l{\cel{3}o per gas-motoro). - DEFIRS.}
\l{}
\l{\sou{elektrono}. (\sou{Fiziko}). Korpuskulo qua portas la maxim mikra karg-}
\l{\cel{3}ajo izolebla ek elektro generale negativa. - DEFIRS.}
}
